% from ega4-8-fr.tex
\oldpage[IV-3]{34}
dans un autre~; on en conclut que d'une façon générale la formation des
$S$-préschémas de présentation finie par la donnée des
$S_\lambda$-préschémas de présentation finie est «~\emph{compatible avec les
limites projectives finies}~», une telle limite étant par définition le noyau
d'un système fini de morphismes d'un même $S$-préschéma dans un produit de
$S$-préschémas (T, 1.8).

Nous nous permettrons couramment par la suite de faire les traductions
impliquées par les propriétés précédentes (ainsi que par (8.3.11), (8.5.2) et
(8.6.3)) sans toujours nous astreindre à les justifier pas à pas comme
ci-dessus. Par exemple, la donnée d'un «~\emph{préschéma en groupes}~» $G$ de
présentation finie sur $S$ équivaut à la donnée d'un préschéma en groupes
$G_\lambda$ de présentation finie sur un $S_\lambda$ pour un $\lambda$ assez
grand~; écrire en effet la condition d'associativité pour le $S$-morphisme
«~loi de composition~» $G\times_SG\to G$ du préschéma en groupes revient à
écrire que le noyau de deux $S$-morphismes de la forme
$G\times_SG\times_SG\to G$ est égal à $G\times_SG\times_SG$ (II, 8.3.9), et
les autres conditions intervenant dans la définition d'un préschéma en
groupes s'interprètent de même.
\end{env}

\subsection{Premières applications à l'élimination des hypothèses noethériennes}
\label{subsection:IV.8.9-fr}

\begin{proposition}[8.9.1]
\EGAReaderTarget{ega4.c04.prop.8-9-1}{IV.8.9.1-fr}
Soient $A$ un anneau, $X$ un $A$-préschéma.
\begin{enumerate}
\item[(i)] Les conditions suivantes sont équivalentes~:
\begin{enumerate}
\item[\emph{a)}] $X$ est de présentation finie sur $A$.
\item[\emph{b)}] Il existe un anneau noethérien $A_0$, un préschéma $X_0$ de
type fini sur $A_0$, un homomorphisme d'anneaux $A_0\to A$, et un
$A$-isomorphisme $X_0\otimes_{A_0}A\xrightarrow{\sim}X$.
\item[\emph{c)}] Il existe un sous-anneau $A_0$ de $A$, qui est une
$\mathbf Z$-algèbre de type fini, un préschéma $X_0$ de type fini sur $A_0$ et
un $A$-isomorphisme $X_0\otimes_{A_0}A\xrightarrow{\sim}X$.
\end{enumerate}
\item[(ii)] Si $\mathcal F$ est un $\mathcal O_X$-Module quasi-cohérent de
présentation finie, on peut supposer le sous-anneau $A_0$ tel qu'il existe un
$\mathcal O_{X_0}$-Module cohérent $\mathcal F_0$ tel que $\mathcal F$ soit
isomorphe à $\mathcal F_0\otimes_{A_0}A$~;
$\operatorname{Supp}(\mathcal F)$ est constructible et fermé dans $X$, et il
y a un sous-préschéma fermé $Z$ de $X$, ayant
$\operatorname{Supp}(\mathcal F)$ comme espace sous-jacent, tel que
l'immersion canonique $Z\to X$ soit de présentation finie.
\item[(iii)] Si $Y$ est un second $A$-préschéma de présentation finie, et
$f:X\to Y$ un $A$-morphisme, on peut supposer le sous-anneau $A_0$ de $A$ tel
qu'il existe un préschéma $Y_0$ de type fini sur $A_0$, un $A$-isomorphisme
$Y_0\otimes_{A_0}A\xrightarrow{\sim}Y$ et un $A_0$-morphisme
$f_0:X_0\to Y_0$ (nécessairement de type fini) tels que $f$ s'identifie à
$f_0\otimes1_A$.
\end{enumerate}
\end{proposition}

\begin{proof}
(i) Comme $A$ est limite inductive de ses sous-anneaux qui sont de type fini
sur $\mathbf Z$, \emph{a)} implique \emph{c)} en vertu de (8.8.2, (ii))~;
\emph{c)} implique \emph{b)} car une $\mathbf Z$-algèbre de type fini est un
anneau noethérien~; enfin, si $A_0$ est noethérien, un $A_0$-préschéma de type
fini est de présentation finie (I, 1.6.1), donc \emph{b)} implique \emph{a)}
en vertu de (I, 1.6.2, (iii)).

L'assertion (ii) se déduit immédiatement de (8.5.2, (ii)), (8.3.11) et
(I, 1.8.2) et l'assertion (iii) de (8.8.2, (i)).
\end{proof}

\begin{env}[8.9.2]
\label{IV.8.9.2-fr}
La proposition (8.9.1) et les résultats de (8.5.2) et (8.8.2) permettent
d'étendre dans de nombreux cas à des morphismes de présentation finie
$X\to Y$ des résultats prouvés par des techniques noethériennes en supposant
$Y$ \emph{localement noethérien}.
